\section{Effective divisors and finite component sums}

Effective divisors on a complex manifold are represented by their invertible
ideal sheaves. The following arithmetic uses these actual ideals, so that a
component factorization yields an equality of effective divisors, including
multiplicities, rather than only an equality of supports.

\begin{definition}[Addition and zero for effective divisors]
\label{hir:effective-divisor-addition}
The zero divisor has ideal $\mathcal O_X$. For effective divisors $D$ and $G$,
define $D+G$ by
$\mathcal O_X(-D-G)=\mathcal O_X(-D)\mathcal O_X(-G)$.
\lean{Hironaka.HironakaEffectiveDivisorSum.instZeroEffectiveDivisor,Hironaka.HironakaEffectiveDivisorSum.instAddEffectiveDivisor}
\leanok
\end{definition}
The unit ideal is invertible, with generator $1$. Products of invertible
ideals are invertible: on the intersection of their generator neighborhoods
the product generator is the product of the two generators. Its nonvanishing
on every nonempty open subset follows by successively avoiding the zero sets
of the two generators. Thus both operations produce effective divisors.

\begin{lemma}[Commutative additive arithmetic]
\label{hir:effective-divisor-additive-monoid}
\uses{hir:effective-divisor-addition}
Effective divisors form a commutative additive monoid. Two effective divisors
are equal if and only if their invertible ideals are equal.
\lean{Hironaka.HironakaEffectiveDivisorSum.instAddCommMonoidEffectiveDivisor,Hironaka.HironakaEffectiveDivisorSum.effectiveDivisor_eq_of_ideal_eq}
\leanok
\end{lemma}
\begin{proof}
The representation of a divisor consists of its ideal and the proposition
that the ideal is invertible, so equality of the ideals gives equality of
the divisors. Commutativity of products of ideals follows by interchanging
the two factors of each product generator. Associativity follows by
expanding finite holomorphic combinations twice and collecting their
products. The unit ideal is neutral because an ideal absorbs holomorphic
combinations of its members, while every member is its own product with
$1$. Invertible ideals are coherent, so these identities apply to all
effective divisors and prove the monoid laws.
\end{proof}

\begin{lemma}[Natural multiples and their support]
\label{hir:effective-divisor-multiples}
\uses{hir:effective-divisor-additive-monoid}
For every $a\in\mathbb N$ and every effective divisor $D$,
$\mathcal O_X(-aD)=\mathcal O_X(-D)^a$.
The support of $0$ is empty, the support of $D+G$ is
$\operatorname{Supp}D\cup\operatorname{Supp}G$, and for $a>0$ the support
of $aD$ is the support of $D$.
\lean{Hironaka.HironakaEffectiveDivisorSum.ideal_nsmul,Hironaka.HironakaEffectiveDivisorSum.support_zero,Hironaka.HironakaEffectiveDivisorSum.support_add,Hironaka.HironakaEffectiveDivisorSum.support_nsmul}
\leanok
\end{lemma}
\begin{proof}
The ideal identity follows by induction on $a$, using the unit ideal at
$a=0$ and multiplication by $\mathcal O_X(-D)$ at the successor step.
At a point outside the zero loci of two ideals choose a nonzero value from
each; their product is a product generator with nonzero value. At a point
of either zero locus every product generator, and hence every holomorphic
combination of them, vanishes. This proves the union identity for support.
The unit ideal contains $1$, so has empty zero locus. The support of a
positive multiple follows by induction and the union identity.
\end{proof}

\begin{lemma}[A smooth hypersurface defines an effective divisor]
\label{hir:smooth-hypersurface-effective-divisor}
A closed smooth hypersurface $H$ defines an effective divisor $[H]$ with
ideal $\mathcal I_H$ and support $H$; the empty hypersurface is permitted.
\lean{Hironaka.HironakaEffectiveDivisorSum.hypersurfaceDivisor,Hironaka.HironakaEffectiveDivisorSum.ideal_hypersurfaceDivisor,Hironaka.HironakaEffectiveDivisorSum.support_hypersurfaceDivisor}
\leanok
\end{lemma}
\begin{proof}
Flattening coordinates identify $H$ locally with the vanishing of one
coordinate. Hadamard division shows that this coordinate generates
$\mathcal I_H$ there. Off the closed set $H$ the generator is $1$.
A hypersurface has empty ambient interior, since a regular point of an
open subset of the ambient manifold has dimension $n$, whereas a regular
point of $H$ has dimension $n-1$. Thus its local generator is identically
zero on no nonempty open subset. This proves invertibility. Members of
$\mathcal I_H$ vanish at every point of $H$, and off $H$ the germ of $1$
is a member; hence its support is exactly $H$.
\end{proof}

\begin{theorem}[Actual finite component sum]
\label{hir:component-product-effective-sum}
\uses{hir:effective-divisor-multiples,hir:smooth-hypersurface-effective-divisor}
For a finite list $A=((H_1,a_1),\ldots,(H_s,a_s))$ of closed smooth
hypersurfaces with natural coefficients, the effective divisor
$\sum_{i=1}^s a_i[H_i]$ has ideal
$\mathcal I_{H_1}^{a_1}\cdots\mathcal I_{H_s}^{a_s}$.
Consequently an equality of an effective divisor's ideal with this product
gives equality of that divisor with this finite sum.
\lean{Hironaka.HironakaEffectiveDivisorSum.ideal_componentSum,Hironaka.HironakaEffectiveDivisorSum.eq_componentSum_of_ideal_eq}
\leanok
\end{theorem}
\begin{proof}
Construct each $[H_i]$ by the preceding lemma. Starting with the zero
divisor, form the finite sum using addition and natural multiplication.
Induction on the list, with the power identity for each summand and the
product identity for addition, gives the asserted ideal. Equality of
ideals then gives equality of effective divisors. This argument permits
the empty list, zero coefficients and repeated hypersurfaces.
\end{proof}
